\documentclass[11pt]{article}

% --- Packages ---
\usepackage[margin=1in]{geometry}
\usepackage{amsmath,amssymb,amsthm}
\usepackage{algorithm}
\usepackage{algorithmic}
\usepackage{booktabs}
\usepackage{graphicx}
\usepackage{hyperref}
\usepackage{natbib}
\usepackage{xcolor}
\usepackage{enumitem}

% --- Theorem environments ---
\newtheorem{theorem}{Theorem}[section]
\newtheorem{lemma}[theorem]{Lemma}
\newtheorem{corollary}[theorem]{Corollary}
\newtheorem{proposition}[theorem]{Proposition}
\newtheorem{definition}[theorem]{Definition}
\newtheorem{remark}[theorem]{Remark}
\newtheorem*{theorem*}{Theorem}
\newtheorem*{corollary*}{Corollary}

% --- Macros ---
\newcommand{\EV}{\mathrm{EV}}
\newcommand{\sub}{\mathrm{sub}}
\newcommand{\cost}{\mathrm{cost}}
\newcommand{\rep}{\mathrm{rep}}
\newcommand{\R}{\mathbb{R}}
\newcommand{\calT}{\mathcal{T}}
\newcommand{\calC}{\mathcal{C}}
\newcommand{\calE}{\mathcal{E}}
\newcommand{\calX}{\mathcal{X}}
\DeclareMathOperator*{\argmin}{arg\,min}

\title{Recursive Bipartite Matching Distance for Game Tree Abstraction\\with Bounded Error}

\author{
  Carmelo Piccione\thanks{Work in progress. Code: \url{https://github.com/struktured-labs/recursive-bipartite-matching}}
}

\date{}

\begin{document}

\maketitle

% ====================================================================
% ABSTRACT
% ====================================================================
\begin{abstract}
We introduce the \emph{recursive bipartite matching} (RBM) distance, a
metric on rooted trees defined by recursively solving minimum-cost
bipartite matchings over children. Unlike scalar hand-strength metrics
used in prior game abstraction work, RBM captures the full strategic
structure of continuation trees. We prove three results: (1)~RBM is a
metric when the leaf distance is a metric; (2)~merging two trees at
distance~$\varepsilon$ introduces expected-value error at most
$\varepsilon/2$ under chance-uniform leaf averaging (a leaf-mean
bound, not yet a Nash exploitability bound; we discuss the bridge to
exploitability via \citet{kroer2014extensive,kroer2018unified} in
Section~\ref{sec:merge}); and (3)~an online learner that clusters game
states by RBM distance achieves cumulative regret
$O(K V_{\max} + T\varepsilon/2)$ with fixed threshold, or $O(\sqrt{T})$
under $\varepsilon$-annealing. On Rhode Island Hold'em, RBM-based
abstraction has lower EV error than Earth Mover's Distance abstraction
across all non-trivial compression levels we evaluate. On full
2-player Limit Hold'em with a standard 52-card deck---the benchmark
game solved by \citet{bowling2015solving}---RBM has lower preflop EV
error than EMD across all five compression levels we measure, with up
to 42\% less error. In a 20{,}000-hand head-to-head match between
MCCFR bots trained on the two abstractions, the RBM bot finishes
+0.02 bb/hand ahead, but the result is not statistically significant
(95\% CI $\approx [-0.19, +0.23]$). On No-Limit Hold'em, RBM wins
abstraction quality at all five compression levels at 20bb (up to 73\%
less error). At 200bb, where pairwise RBM distance is computed on
showdown-distribution trees rather than full game trees, RBM dominates
fine-grained compression (60\% less error at $k{=}25$, 40\% at
$k{=}15$) but loses at coarser levels.
\end{abstract}

% ====================================================================
% 1. INTRODUCTION
% ====================================================================
\section{Introduction}
\label{sec:intro}

Abstraction is the dominant paradigm for solving large extensive-form
games. The idea is straightforward: collapse strategically similar game
states into equivalence classes, solve the smaller \emph{abstract} game,
and map the abstract strategy back to the original.  Superhuman poker
agents~\citep{bowling2015solving,moravvcik2017deepstack,brown2018superhuman}
all rely on some form of abstraction to make computation tractable.

The quality of an abstraction depends critically on the \emph{distance
metric} used to decide which states are ``similar.'' The standard
approach in poker AI, pioneered by \citet{gilpin2007potential}, clusters
information sets using Earth Mover's Distance (EMD) on hand-strength
histograms: distributions over scalar equity values obtained by
evaluating each hand against a uniform random opponent. This is
effective but inherently limited---it reduces each game state to a
one-dimensional summary, discarding the branching structure of future
play.

Two hands may have identical equity distributions yet lead to
fundamentally different strategic situations. A made hand and a draw can
have the same expected value but completely different continuation trees:
one calls safely, the other bets aggressively to deny correct pot odds
or folds against reraises. Scalar metrics cannot distinguish them.

We propose a structural alternative. The \emph{recursive bipartite
matching} (RBM) distance between two game trees is defined by
recursively computing minimum-cost bipartite matchings between their
children at every level. Leaves are compared by payoff difference;
internal nodes are compared by the cost of optimally aligning their
subtrees. Structural mismatches (differing numbers of children) incur an
explicit phantom penalty. The result is a metric that respects both the
values at the leaves and the \emph{shape} of the computation that
produces them.

This paper develops the theory and practice of RBM-based game
abstraction in three parts.

\paragraph{The metric and its error bound.}  We prove that RBM is a
metric (Section~\ref{sec:rbm}), and that merging two trees at
distance~$\varepsilon$ introduces EV error at most~$\varepsilon/2$
(Section~\ref{sec:merge}). The bound is tight at leaves but strictly
tighter at branching nodes, where averaging reduces error.

\paragraph{The location problem.}  Using a structurally rich distance
creates a compressed \emph{EV graph}---a DAG of merged game
states---that faithfully preserves strategic content. But locating
oneself in this graph during live play requires comparing one's
continuation tree against cluster representatives, which can be as
expensive as solving the uncompressed game
(Section~\ref{sec:location}).

\paragraph{The online learning resolution.}  We dissolve the location
problem by building the EV graph \emph{incrementally through self-play}.
The learner always knows where it is because it constructed the graph
from its own positions. We prove that this online learner achieves
bounded cumulative regret, with sublinear average regret under
$\varepsilon$-annealing (Section~\ref{sec:online}).

Experiments on Rhode Island Hold'em~\citep{gilpin2005rhode}, full
2-player Limit Hold'em~\citep{bowling2015solving}, and No-Limit
Hold'em~\citep{brown2018superhuman} (Section~\ref{sec:experiments})
report consistent improvements in abstraction quality: RBM-based
clustering produces lower EV error than EMD across all non-trivial
compression levels we evaluate, on Rhode Island Hold'em, on Limit
Hold'em (5--0 at $k \in \{3,5,10,15,25\}$, up to 42\% less error), and
on No-Limit 20bb (up to 73\% less error at $k{=}25$). At 200bb, RBM
dominates only at fine-grained compression ($k \ge 15$) and loses at
coarser levels---an asymmetry consistent with the
abstraction-pathology phenomenon documented by
\citet{waugh2009abstraction} and amplified here because the 200bb
experiment uses showdown-distribution trees as a proxy for the full
186{,}174-node game tree (full-tree RBM at 200bb remains future work,
Section~\ref{sec:conclusion}). MCCFR bots trained on RBM versus EMD
abstractions in Limit Hold'em finish within statistical noise of each
other in head-to-head play (RBM $+$0.02 bb/h, $n{=}20{,}000$, 95\% CI
$\approx [-0.19, +0.23]$): no detectable gameplay difference, despite
the abstraction-quality gap. The online learner achieves 99.4\% cache
hit rates with ${\sim}400$ naturally emerging postflop clusters.

% ====================================================================
% 2. RECURSIVE BIPARTITE MATCHING DISTANCE
% ====================================================================
\section{Recursive Bipartite Matching Distance}
\label{sec:rbm}

\subsection{Definition}

Let $\calT$ denote the space of finite rooted trees with real-valued
leaves. We define a distance $d : \calT \times \calT \to \R_{\ge 0}$
recursively.

\begin{definition}[RBM Distance]
\label{def:rbm}
Given two rooted trees $T_1$ and $T_2$:

\textbf{Base case.} If both are leaves with values $v_1, v_2$:
\[
  d(\mathrm{leaf}_1, \mathrm{leaf}_2) = |v_1 - v_2|.
\]

\textbf{Recursive case.} Let $n_1$ have children $\{c_1^1, \ldots,
c_1^m\}$ and $n_2$ have children $\{c_2^1, \ldots, c_2^n\}$, with $m
\le n$ (WLOG).
\begin{enumerate}[nosep]
  \item Compute the pairwise cost matrix $W$ where $W[i,j] =
    d(\sub(c_1^i), \sub(c_2^j))$.
  \item Pad the smaller child set with $|n - m|$ \emph{phantom nodes}.
    The cost of matching a real subtree $S$ against a phantom is
    $\delta(S)$.
  \item Solve the minimum-cost perfect matching on the padded bipartite
    graph via the Hungarian method~\citep{kuhn1955hungarian}, yielding
    optimal matching $M^*$.
  \item $d(T_1, T_2) = \cost(M^*)$.
\end{enumerate}
\end{definition}

The phantom penalty $\delta$ quantifies the cost of structural
mismatch. The cleanest formulation is $\delta(S) = d(S, \emptyset)$,
treating the phantom as a designated empty tree in the metric space.
Under this choice, all conditions needed for the metric proof are
consequences of the triangle inequality at lower depth (see
Remark~\ref{rem:phantom}). To make this rigorous we extend the base
case of Definition~\ref{def:rbm} to cover $\emptyset$ explicitly:
$d(\mathrm{leaf}(v), \emptyset) = |v|$ and $d(T, \emptyset) =
\sum_{c \in \mathrm{children}(T)} d(\sub(c), \emptyset)$ when $T$ is
internal. With this completion, $\delta$ is well-defined for every
tree, and Remark~\ref{rem:phantom} below recovers the 1-Lipschitz and
phantom-triangle properties as instances of the lower-depth triangle
inequality.

Children are treated as \emph{unordered sets}. The matching finds the
best structural alignment without requiring a canonical ordering---exactly
the right semantics for game trees, where children of a chance node
(possible deals) or decision node (possible actions) have no inherent
order.

\subsection{Metric Properties}

\begin{theorem}[RBM is a Metric]
\label{thm:metric}
Let $d_L$ be a metric on leaf values and let $\delta(S) = d(S,
\emptyset)$ for a designated empty tree $\emptyset$. Then the RBM
distance $d$ is a metric on~$\calT$.
\end{theorem}

\begin{proof}[Proof sketch]
By structural induction on tree depth. The base case ($k=0$) reduces to
$d_L$, which is a metric by assumption.

\emph{Identity:} The identity matching $\{(c_i, c_i)\}$ has cost zero
by induction; since all costs are nonneg, the minimum-cost matching
achieves $d(T, T) = 0$.

\emph{Symmetry:} The cost matrix for $d(T_2, T_1)$ is the transpose of
that for $d(T_1, T_2)$, using symmetry at depth ${<}\,k$. The
minimum-cost perfect matching is invariant under transposition.

\emph{Triangle inequality:} Given optimal matchings $M_{12}$ (between
$T_1, T_2$) and $M_{23}$ (between $T_2, T_3$), compose through $T_2$
to construct a feasible matching $M_{13}$ between $T_1$ and $T_3$.
Children ``through-matched'' via some $b_j \in T_2$ satisfy
$d(\sub(a_i), \sub(c_k)) \le d(\sub(a_i), \sub(b_j)) + d(\sub(b_j),
\sub(c_k))$ by induction. Dangling children (matched to phantoms on one
side) are handled by the 1-Lipschitz property of~$\delta$. Summing over
all pairs: $\cost(M_{13}) \le \cost(M_{12}) + \cost(M_{23})$. Since
$d(T_1, T_3) \le \cost(M_{13})$ by optimality, the triangle inequality
follows.

The full proof appears in Appendix~\ref{app:metric}.
\end{proof}

\begin{remark}[Phantom Penalty]
\label{rem:phantom}
With $\delta(S) = d(S, \emptyset)$, the phantom triangle property
$d(S_1, S_2) \le \delta(S_1) + \delta(S_2)$ and the 1-Lipschitz
condition $\delta(S_1) \le d(S_1, S_2) + \delta(S_2)$ are both
instances of the triangle inequality. The theorem is self-contained: if
the leaf distance is a metric, then RBM with empty-tree phantoms is a
metric.
\end{remark}

\subsection{Relation to Prior Tree Distances}

RBM is related to but distinct from tree edit distance
\citep{zhang1996constrained}, which uses insert/delete/relabel
operations on ordered trees. Edit distance is more flexible but less
natural for \emph{unordered} trees where children are interchangeable.
The recursive matching structure also connects to hierarchical optimal
transport~\citep{memoli2011gromov} and the Tree Mover's
Distance~\citep{chuang2022tree}, but those works target different
applications (shape matching and GNN analysis, respectively). See
Section~\ref{sec:related} for a detailed comparison.

% ====================================================================
% 3. GAME TREE COMPRESSION VIA MERGE
% ====================================================================
\section{Game Tree Compression via Merge}
\label{sec:merge}

\subsection{The Merge Operation}

The RBM distance induces a natural \emph{merge} operation. Given trees
$T_1, T_2$ with small distance, produce a representative $T^*$ that
captures their shared structure.

\begin{definition}[Tree Merge]
Given the optimal matching $M^*$ between children of $T_1$ and $T_2$:
\begin{itemize}[nosep]
  \item \textbf{Matched pairs} $(c_1^i, c_2^j) \in M^*$: recursively
    merge their subtrees.
  \item \textbf{Phantom-matched children}: either retain at reduced
    weight or drop (lossy compression).
  \item \textbf{Leaf merging}: average the payoff values: $v^* =
    (v_1 + v_2) / 2$.
\end{itemize}
\end{definition}

\subsection{The EV Graph}

Repeated merging across all subtrees at each depth produces a
directed acyclic graph---the \emph{EV graph}. Multiple original game
paths converge to the same node, meaning ``from here, these situations
play out essentially the same way.'' This is hierarchical clustering
over subtrees using RBM as the metric, with each cluster represented by
a merged prototype tree.

\subsection{EV Error Bound}

\begin{theorem}[EV Error Bound Under Merging]
\label{thm:ev-error}
Let $T_1, T_2$ be rooted trees with real-valued leaves, $d(T_1, T_2) =
\varepsilon$. Let $T^* = \mathrm{merge}(T_1, T_2)$ with equal-weight
leaf averaging. Then:
\[
  |\EV(T^*) - \EV(T_i)| \le \frac{\varepsilon}{2}, \quad i \in \{1, 2\}.
\]
\end{theorem}

\begin{proof}[Proof sketch]
By structural induction on tree depth.

\emph{Base case (leaves):} $T^* = \mathrm{leaf}((v_1+v_2)/2)$, so
$|\EV(T^*) - v_1| = |v_2 - v_1|/2 = \varepsilon/2$.

\emph{Inductive case:} Let $M^*$ be the optimal matching with
matched-pair costs $d_i$ and total $\varepsilon = \sum_i d_i +
\sum_{\text{phantom}} \delta(b_j)$. Each merged child satisfies
$|\EV(c_i^*) - \EV(\sub(a_i))| \le d_i / 2$ by induction. With uniform
weighting over $m$ children:
\[
  |\EV(T^*) - \EV(T_1)| \le \frac{1}{m} \sum_{i=1}^m \frac{d_i}{2}
  = \frac{1}{2m} \sum_{i=1}^m d_i \le \frac{\varepsilon}{2}.
\]

For trees with branching factor $m > 1$ and no phantom terms, the bound
tightens to $\varepsilon / (2m) \le \varepsilon / 4$. The stated bound
$\varepsilon/2$ is tight only at leaves or degenerate single-child
chains.
\end{proof}

\begin{remark}[EV bound vs.\ exploitability]
\label{rem:ev-vs-expl}
Theorem~\ref{thm:ev-error} bounds $|\EV(T^*) - \EV(T_i)|$, where $\EV$
is the expectation over the tree's leaves under the chance distribution
already baked into the leaf weights---the EV that a fixed strategy
profile delivers when one substitutes the merged tree for the original.
This is \emph{not} the same as a bound on Nash exploitability, which
asks how much a best-responder can win against a strategy obtained by
solving the abstract game and lifting back to the full game. The
canonical exploitability bounds for game-tree abstractions are due to
\citet{kroer2014extensive,kroer2018unified}; a tight bridge from
RBM-induced merge error to exploitability would proceed via their
level-by-level reward-and-action distance machinery, applied to the
specific child-matching structure RBM produces.  We treat that bridge
as the natural follow-up; the present bound should be read as a
necessary but not sufficient condition for low exploitability.
\end{remark}

\begin{corollary}[Weighted Merges]
For unequal-weight merges with weights $w_1 + w_2 = 1$:
\[
  |\EV(T^*) - \EV(T_i)| \le w_{3-i} \cdot \varepsilon.
\]
Equal weights $w_1 = w_2 = 1/2$ recover $\varepsilon/2$.
\end{corollary}

\begin{corollary}[Iterated Merges]
When $k$ trees $T_1, \ldots, T_k$ are merged into a single $T^*$ via
successive pairwise merges:
\[
  |\EV(T^*) - \EV(T_i)| \le \max_j d(T_i, T_j) \le \varepsilon_{\mathrm{diam}},
\]
where $\varepsilon_{\mathrm{diam}}$ is the cluster diameter.
\end{corollary}

% ====================================================================
% 4. THE LOCATION PROBLEM
% ====================================================================
\section{The Location Problem}
\label{sec:location}

Building the EV graph offline is one challenge; \emph{finding yourself
in it during live play} is another.

\subsection{Statement}

Given a compressed EV graph $G$ and a game state $v$ (known: cards,
history), find the corresponding node $\phi(v) \in G$.

In flat abstraction, location is trivial: precompute a lookup table
mapping each state to its bucket, $O(1)$ at query time. With tree-based
abstraction, the equivalence class of $v$ depends on its \emph{entire
continuation subtree}---the same non-local structure that makes the
distance expressive.

\subsection{Complexity Barrier}

\begin{proposition}[Location Cost]
\label{prop:location}
Let $b$ be the branching factor, $h$ the remaining depth, $K$ the
number of clusters. Naive location requires $O(K \cdot b^{2h})$ time
per game step, since each distance computation $T(h)$ satisfies $T(h) =
b^2 \cdot T(h-1) + O(b^3)$, giving $T(h) = O(b^{2h})$.
\end{proposition}

Over a full game of $d$ steps, the total cost is $O(d \cdot K \cdot
n^2)$ where $n = b^h$ is the number of leaves. This can exceed the cost
of solving the uncompressed game, defeating the purpose of compression.
The location problem likely admits a lower bound of $\Omega(n)$ per
step, because the equivalence classes depend on global subtree
structure defined by optimal (Hungarian) matchings.

This reveals a fundamental tradeoff in game abstraction: \emph{the more
structurally faithful the abstraction, the harder it is to navigate at
runtime.}

% ====================================================================
% 5. ONLINE LEARNING WITH BOUNDED REGRET
% ====================================================================
\section{Online Learning with Bounded Regret}
\label{sec:online}

The offline approach---enumerate, cluster, then look up---hits the
location problem wall. We propose a more natural formulation: build the
EV graph \emph{dynamically through play}.

\subsection{The Online EV Graph Learner}

\begin{algorithm}[t]
\caption{Online EV Graph Construction via Self-Play}
\label{alg:online}
\begin{algorithmic}[1]
\STATE Initialize EV graph $G \leftarrow \emptyset$, threshold
  $\varepsilon > 0$
\FOR{each game $t = 1, 2, \ldots, T$}
  \FOR{each decision point in game $t$}
    \STATE Observe game state $s$; build partial subtree $T_s$
    \STATE Find nearest cluster: $c^* = \argmin_c d(T_s, \rep(c))$
    \IF{$d(T_s, \rep(c^*)) \le \varepsilon$ \textbf{(cache hit)}}
      \STATE Use cluster $c^*$'s strategy; record outcome
    \ELSE[\textbf{(cache miss)}]
      \STATE Create new cluster from $T_s$; explore
    \ENDIF
  \ENDFOR
  \STATE Update cluster representatives with observed payoffs
  \STATE Merge clusters where $d(\rep(c_i), \rep(c_j)) <
    \varepsilon_{\mathrm{merge}}$
  \STATE Optionally anneal $\varepsilon$ downward
\ENDFOR
\end{algorithmic}
\end{algorithm}

Algorithm~\ref{alg:online} resolves the core tensions:

\paragraph{The location problem dissolves.} The learner never needs to
locate itself in a precomputed graph because it built the graph
\emph{from} its own positions. Each cluster knows its game states
because it was born from one.

\paragraph{The enumeration problem dissolves.} Only states actually
encountered are represented. The graph concentrates resolution where it
matters---the same insight behind MCTS~\citep{coulom2006efficient}.

\paragraph{The error bound becomes a regret bound.}  When the learner
uses cluster $c^*$'s strategy on a state at distance $\le \varepsilon$,
the EV loss is bounded by $\varepsilon/2$ (Theorem~\ref{thm:ev-error}).

\subsection{Regret Bound}

\begin{theorem}[Online Regret Bound]
\label{thm:regret}
Consider the online EV graph learner (Algorithm~\ref{alg:online}) with
threshold $\varepsilon > 0$, playing $T$ games against a stationary
environment. Let $K$ be the number of distinct strategic equivalence
classes (game trees up to $\varepsilon$-distance), and $V_{\max} =
\max_t |\EV(T_t)|$. Then:

\begin{enumerate}[nosep]
  \item \textbf{Exploration bound:} Total cache misses $\le K$.
  \item \textbf{Per-step exploitation error:} At most $\varepsilon/2$
    per cache hit.
  \item \textbf{Cumulative regret:}
    \[
      R(T) \le K \cdot V_{\max} + (T - K) \cdot
      \frac{\varepsilon}{2}.
    \]
  \item \textbf{Average regret:} $R(T)/T \to \varepsilon/2$ as $T \to
    \infty$.
  \item \textbf{$\varepsilon$-annealing:} With $\varepsilon(t) =
    \varepsilon_0 / \sqrt{t}$, $R(T) = O(\sqrt{T})$ and $R(T)/T \to 0$.
\end{enumerate}
\end{theorem}

\begin{proof}[Proof sketch]
\emph{Part~1:} Each equivalence class triggers at most one cache miss.
Once representative $T_t$ is created, all future trees in the same class
satisfy $d(T_{t'}, T_t) \le \varepsilon$, producing cache hits.

\emph{Parts~2--3:} Cache hits incur error $\le \varepsilon/2$
(Theorem~\ref{thm:ev-error}). Cache misses incur at most $V_{\max}$.
Summing: $R(T) \le K V_{\max} + (T-K)\varepsilon/2$.

\emph{Part~5:} With $\varepsilon(t) = \varepsilon_0/\sqrt{t}$,
exploitation regret sums to $\sum_{t=1}^T \varepsilon_0/(2\sqrt{t}) \le
\varepsilon_0 \sqrt{T}$. Exploration cost is at most $N \cdot V_{\max}$
(constant). Total: $R(T) = O(\sqrt{T})$.

The full proof appears in Appendix~\ref{app:regret}.
\end{proof}

\subsection{Connection to Lipschitz Bandits}

The $O(\sqrt{T})$ rate matches the minimax-optimal rate for multi-armed
bandits~\citep{auer2002finite}. The structural advantage over generic
bandits is that RBM provides a \emph{geometry} on the arm space: nearby
clusters have similar payoffs by the EV error bound. This is analogous
to Lipschitz bandits or continuum-armed bandits~\citep{kleinberg2005nearly},
where metric structure enables faster learning. In the EV graph setting,
the metric is not assumed but \emph{derived} from the game structure via
recursive bipartite matching.

% ====================================================================
% 6. EXPERIMENTS
% ====================================================================
\section{Experiments}
\label{sec:experiments}

We evaluate on three domains: (1)~Rhode Island
Hold'em~\citep{gilpin2005rhode}, a simplified poker variant (2 players,
1 hole card, 2 community cards, 3 betting rounds, limit betting),
(2)~full 2-player Limit Hold'em with a standard 52-card deck---the
benchmark game solved by \citet{bowling2015solving}, and (3)~No-Limit
Hold'em (heads-up, 20bb stacks)---the game variant underlying
modern superhuman agents~\citep{brown2018superhuman}.

The implementation is in OCaml (${\sim}$15K lines) using Jane Street
Core libraries, with an $O(n^3)$ Hungarian algorithm for bipartite
matching.

\subsection{RBM vs EMD Comparison}

We compare three distance metrics for game abstraction:
\begin{itemize}[nosep]
  \item \textbf{RBM}: recursive bipartite matching on full game trees
    ($\sim$1141 nodes per deal),
  \item \textbf{EMD}: Earth Mover's Distance on showdown outcome
    distributions (win/draw/lose probabilities), the standard poker AI
    metric~\citep{gilpin2007potential},
  \item \textbf{Scalar EV}: absolute difference in expected values.
\end{itemize}

For each metric, we compute pairwise distances on 25 sampled deals
(3-rank deck, fixed P1 card, varied opponent/community), then perform
agglomerative clustering at multiple $\varepsilon$ thresholds. At each
compression level (cluster count), we measure the maximum EV error
across clusters.

\begin{table}[t]
\centering
\caption{RBM vs EMD vs Scalar EV: maximum EV error at matched cluster
counts (3-rank Rhode Island Hold'em, 25 sampled deals with varied
community cards). Lower is better. RBM achieves zero error down to
$k=3$ clusters while EMD has 30.28 error at every $k \ge 3$. The
constant 30.28 reflects a degeneracy of EMD on this 3-rank toy: the
showdown-distribution histograms collapse to a small set of equity
values, so refining the cluster count cannot reduce error past a
floor. We therefore treat this table as illustrative of the
representational ceiling of scalar metrics on toy decks, not as a
benchmark; the more informative comparisons are the 52-card
abstraction-quality results in Sections~\ref{sec:experiments-limit}
and~\ref{sec:experiments-nl}.}
\label{tab:rbm-vs-emd}
\begin{tabular}{rccccl}
\toprule
Clusters & Compress & RBM Error & EMD Error & EV Error & Winner \\
\midrule
25 & 1.0$\times$ & 0.00 & 30.28 & 30.28 & RBM \\
15 & 1.7$\times$ & 0.00 & 30.28 & 30.28 & RBM \\
10 & 2.5$\times$ & 0.00 & 30.28 & 30.28 & RBM \\
 7 & 3.6$\times$ & 0.00 & 30.28 & 30.28 & RBM \\
 5 & 5.0$\times$ & 0.00 & 30.28 & 30.28 & RBM \\
 3 & 8.3$\times$ & 0.00 & 30.28 & 30.28 & RBM \\
 2 & 12.5$\times$ & 25.96 & 34.07 & 34.07 & RBM \\
 1 & 25.0$\times$ & 54.51 & 54.51 & 54.51 & tie \\
\bottomrule
\end{tabular}
\end{table}

Table~\ref{tab:rbm-vs-emd} shows that on this 3-rank toy RBM produces
zero EV error at every compression level $k \ge 2$, while EMD's error
sits at the 30.28 floor described above. This is mostly a statement
about EMD's failure mode on a deck this small: with so few distinct
showdown histograms, the metric simply cannot resolve hands that have
identical showdown distributions but different game-tree continuations.
RBM, operating on the full game tree, separates these by construction.
We include the table to illustrate the failure mode---not as evidence
about which metric is better at production scale.

\subsection{CFR on Compressed Games}

We run counterfactual regret minimization (CFR)~\citep{zinkevich2007regret}
on both the full game and RBM-compressed abstractions to measure
exploitability---the gap between the computed strategy and Nash
equilibrium.

\begin{table}[t]
\centering
\caption{CFR exploitability on full vs.\ RBM-compressed Rhode Island
Hold'em (3-rank, 2 betting rounds, 1000 CFR iterations). RBM
compression to 3 clusters achieves \textbf{zero exploitability}---the
abstraction is strategically lossless.}
\label{tab:cfr}
\begin{tabular}{rrrr}
\toprule
$\varepsilon$ & Clusters & EV Error & Exploitability \\
\midrule
(full) & 10 & --- & 0.0000 \\
0 & 3 & 0.0000 & 0.0000 \\
50 & 3 & 0.0000 & 0.0000 \\
100 & 3 & 0.0000 & 0.0000 \\
200 & 3 & 0.0000 & 0.0000 \\
500 & 3 & 0.0000 & 0.0000 \\
1000 & 1 & 7.0926 & 5.6573 \\
\bottomrule
\end{tabular}
\end{table}

Table~\ref{tab:cfr} shows that on this 3-rank, 10-info-set toy, RBM
compression to 3 clusters preserves zero exploitability across the
$\varepsilon$ range we tested. The natural reading is that RBM has
correctly identified the three rank-vs-board equivalence classes the
game admits at this scale; we are reluctant to call the result
``strategically lossless'' in any general sense, because the result
also depends on the unique structure of the 3-rank game (with only 10
information sets, there is little room for an abstraction to do
damage). Only at $\varepsilon = 1000$, where the algorithm collapses
all hands into one cluster, does exploitability rise---consistent with
Theorem~\ref{thm:ev-error}'s leaf-mean bound but not, on its own,
evidence of a tight exploitability bound at larger scales (cf.\
Remark~\ref{rem:ev-vs-expl}).

\subsection{Multi-Scale Results}

Table~\ref{tab:multiscale} shows RBM vs EMD across three toy-deck scales.
RBM achieves lower error across these target compression levels.
We note, however, that these ``scales'' (3-rank, 4-rank, and 5-rank) are
still within the family of highly compressed toy environments, and the
measured dominance over EMD is largely an artifact of using a
potential-blind showdown-only EMD baseline (which lacks the future card-draw
distributions used in production-scale EMD abstractions; cf.\
Section~\ref{sec:boundary}). When EMD is stripped of potential-awareness,
it cannot resolve the intermediate game-tree paths that RBM captures
by construction. We present this comparison to verify that RBM's structural
advantage is consistent as the toy deck size increases, not as a
generalization to full-scale games (which are evaluated in
Sections~\ref{sec:experiments-limit} and~\ref{sec:experiments-nl}).

\begin{table}[t]
\centering
\caption{RBM vs EMD across game scales (30 sampled deals per scale,
varied community cards). RBM wins column counts victories across 8
target compression levels. RBM achieves zero error down to $k=2$
at all scales.}
\label{tab:multiscale}
\begin{tabular}{lccccc}
\toprule
Scale & Cards & RBM wins & EMD wins & Ties & RBM zero-err $k$ \\
\midrule
3-rank & 12 & 7 & 0 & 1 & 2 \\
4-rank & 16 & 7 & 0 & 1 & 2 \\
5-rank & 20 & 7 & 0 & 1 & 2 \\
\midrule
\textbf{Total} & & \textbf{21} & \textbf{0} & \textbf{3} & \\
\bottomrule
\end{tabular}
\end{table}

\subsection{Metric Property Verification}

We verify RBM metric properties exhaustively on the 4-rank game
(91~deal trees, 753,571 triples for the triangle inequality):

\begin{itemize}[nosep]
  \item \textbf{Identity}: $d(T, T) = 0$ for all 91 trees. \checkmark
  \item \textbf{Symmetry}: $d(T_i, T_j) = d(T_j, T_i)$ for all 4095
    pairs. \checkmark
  \item \textbf{Triangle inequality}: verified on all 753,571 triples.
    \checkmark
\end{itemize}

\subsection{Online Learner Convergence}

We run the online self-play learner (Algorithm~\ref{alg:online}) on
2-rank Rhode Island Hold'em with $\varepsilon = 50$ and random community
cards.

\begin{table}[t]
\centering
\caption{Online learner convergence on 2-rank Rhode Island Hold'em.
With random community cards, the learner discovers the natural cluster
structure by game~100 and achieves near-perfect cache hit rates by
game~600.}
\label{tab:online}
\begin{tabular}{rrr}
\toprule
Games & Clusters & Cache Hit Rate \\
\midrule
100 & 3 & 95.0\% \\
200 & 3 & 97.5\% \\
400 & 3 & 98.8\% \\
600 & 3 & 99.5\% \\
\bottomrule
\end{tabular}
\end{table}

Table~\ref{tab:online} shows rapid convergence: the learner discovers
all 3 natural clusters by game~100 (from 100 games with random
community cards), and cache hit rate reaches 99.5\% by game~600. With
fixed community cards, convergence is even faster (same 3 clusters by
game~50), confirming that the smaller strategic space is easier to
learn.

\subsection{Distance Memoization}

The recursive structure of RBM creates many redundant subtree
comparisons. Memoizing subtree distances yields significant speedups.

On an 8-tree benchmark (2-rank deck, 28 pairwise comparisons), memoized
distance computation achieves a \textbf{94$\times$ speedup} over the
naive implementation, with a memo cache hit rate of approximately 97\%.
This demonstrates that game trees share substantial subtree structure
that memoization exploits.

\subsection{Full Limit Hold'em: Preflop Abstraction}
\label{sec:experiments-limit}

We scale to full 2-player Limit Hold'em with a standard 52-card
deck~\citep{bowling2015solving}. After suit isomorphism, there are 50
canonical preflop hands. We compare RBM and EMD abstraction quality
across 5 compression levels by measuring the mean EV error of each
clustering.

\begin{table}[t]
\centering
\caption{RBM vs EMD: preflop abstraction quality on full Limit Hold'em
(50 canonical hands, 52-card deck). Mean EV error at each cluster count.
Lower is better. \textbf{RBM wins all 5 levels.}}
\label{tab:holdem-preflop}
\begin{tabular}{rcccl}
\toprule
Clusters ($k$) & RBM Error & EMD Error & RBM Advantage & Winner \\
\midrule
25 & 0.25 & 0.43 & 42\% less error & RBM \\
15 & 0.37 & 0.43 & 14\% less error & RBM \\
10 & 0.40 & 0.51 & 22\% less error & RBM \\
 5 & 0.60 & 0.69 & 13\% less error & RBM \\
 3 & 0.62 & 0.71 & 13\% less error & RBM \\
\bottomrule
\end{tabular}
\end{table}

Table~\ref{tab:holdem-preflop} shows that \textbf{RBM wins 5--0} across
all compression levels. The advantage is largest at $k = 25$ (42\% less
error), where RBM's structural sensitivity matters most: it correctly
separates hands that have similar equity distributions but different
continuation tree structures (e.g., suited connectors vs.\ offsuit
broadways). Even at extreme compression ($k = 3$), RBM maintains a 13\%
advantage.

\subsection{Full Limit Hold'em: MCCFR Head-to-Head}

We train MCCFR (Monte Carlo Counterfactual Regret Minimization) bots
using RBM-based and EMD-based preflop abstractions, then evaluate in a
20{,}000-hand position-alternated match.

\begin{table}[t]
\centering
\caption{MCCFR bot tournament on full Limit Hold'em (20{,}000 hands,
position-alternated). All trained bots beat both baselines. The RBM bot
beats the EMD bot head-to-head, though the margin is not statistically
significant (see text).}
\label{tab:holdem-cfr}
\begin{tabular}{lc}
\toprule
Matchup & Result (bb/hand) \\
\midrule
RBM bot vs EMD bot & $+0.02$ (RBM wins)\textsuperscript{$\dagger$} \\
RBM bot vs Random & $+1.24$ \\
EMD bot vs Random & $+1.12$ \\
RBM bot vs Always-Call & $+0.30$ \\
EMD bot vs Always-Call & $+0.11$ \\
\bottomrule
\end{tabular}
\end{table}

The head-to-head margin in Table~\ref{tab:holdem-cfr} is $+0.02$
bb/hand---directionally consistent with the abstraction-quality
results in Table~\ref{tab:holdem-preflop}, but well within sampling
noise. We do \emph{not} claim that better abstraction quality
translated to stronger play in this experiment: at $n = 20{,}000$
hands the 95\% CI of approximately $[-0.19, +0.23]$ comfortably covers
zero. The same caveat applies to the differences against the random
and always-call baselines, where both bots win by similar margins
($+1.24$ vs.\ $+1.12$ vs.\ random; $+0.30$ vs.\ $+0.11$ vs.\
always-call) but the differences are consistent with noise. A
larger-$n$ rerun, or a Slumbot-style benchmark match against a fixed
strong opponent (cf.\ \citealt{johanson2013evaluating}), would be the
appropriate follow-up.

\noindent\textsuperscript{$\dagger$}\textit{Statistical note:} The
$+0.02$ bb/hand head-to-head margin is not statistically significant.
With $\sigma \approx 15$ bb/hand for Limit Hold'em, $\mathrm{SE}(20{,}000)
\approx 0.11$ bb/hand, giving a 95\% CI of approximately
$[-0.19, +0.23]$ bb/hand. The directional advantage is consistent with
the abstraction quality results (Table~\ref{tab:holdem-preflop}) but
a larger sample is needed to confirm statistical significance.

\subsection{Full Limit Hold'em: Online Learner}

The online learner (Algorithm~\ref{alg:online}) scales to full Limit
Hold'em, achieving a 99.4\% cache hit rate---60$\times$ faster than
offline distance computation. Approximately 400 postflop clusters emerge
naturally with $\varepsilon = 0.5$, without any hand-crafted bucketing.
This demonstrates that the RBM distance discovers meaningful strategic
groupings even in the full game.

\subsection{No-Limit Hold'em: Preflop Abstraction}
\label{sec:experiments-nl}

We extend the evaluation to heads-up No-Limit Hold'em at two stack
depths: 20bb (short-stack) and 200bb (deep-stack), the game variants
underlying modern superhuman
agents~\citep{brown2018superhuman,moravvcik2017deepstack}. No-Limit
introduces variable bet sizes that create qualitatively different game
trees from Limit poker (Table~\ref{tab:tree-sizes}).

\begin{table}[t]
\centering
\caption{Game tree sizes across poker variants (2-player heads-up unless
noted). No-Limit games are substantially larger due to variable bet
sizing, despite shallower stack depths.}
\label{tab:tree-sizes}
\begin{tabular}{lrl}
\toprule
Variant & Nodes & Notes \\
\midrule
Limit HU & 9{,}476 & Fixed bet sizes \\
NL-HU 20bb & 2{,}988 & Short-stack push/fold \\
NL-HU 200bb & 186{,}174 & Deep-stack play \\
NL 6-max 20bb & ${\sim}$13M & Multi-player \\
\bottomrule
\end{tabular}
\end{table}

\begin{table}[t]
\centering
\caption{RBM vs EMD: preflop abstraction quality on NL-HU 20bb
(50 canonical hands). Max EV error at each cluster count. Lower is
better. \textbf{RBM wins 5--0} across all compression levels.}
\label{tab:nl-preflop}
\begin{tabular}{rcccl}
\toprule
$k$ & RBM Error & EMD Error & Improvement & Winner \\
\midrule
25 & 0.11 & 0.41 & 73\% less error & RBM \\
15 & 0.14 & 0.38 & 63\% less error & RBM \\
10 & 0.25 & 0.38 & 34\% less error & RBM \\
 5 & 0.44 & 0.49 & 10\% less error & RBM \\
 3 & 0.45 & 0.48 & 5\% less error & RBM \\
\bottomrule
\end{tabular}
\end{table}

Table~\ref{tab:nl-preflop} shows the 20bb results. \textbf{RBM wins
5--0} across all compression levels, with the largest advantage at
fine-grained compression (73\% less error at $k{=}25$, 63\% at
$k{=}15$).

\begin{table}[t]
\centering
\caption{RBM vs EMD: preflop abstraction quality on NL-HU 200bb
(50 canonical hands, showdown distribution trees). Max EV error at each
cluster count. Lower is better. RBM dominates at fine-grained
compression; EMD wins at coarse levels.}
\label{tab:nl-200bb-preflop}
\begin{tabular}{rcccl}
\toprule
$k$ & RBM Error & EMD Error & Improvement & Winner \\
\midrule
25 & 0.17 & 0.43 & 60\% less error & RBM \\
15 & 0.28 & 0.46 & 40\% less error & RBM \\
10 & 0.60 & 0.42 & --- & EMD \\
 5 & 0.62 & 0.57 & --- & EMD \\
 3 & 0.62 & 0.58 & --- & EMD \\
\bottomrule
\end{tabular}
\end{table}

Table~\ref{tab:nl-200bb-preflop} shows the 200bb deep-stack results.
RBM dominates at fine-grained compression (60\% less error at
$k{=}25$, 40\% at $k{=}15$) but loses at coarser compression
($k \le 10$). Two effects compound to produce this asymmetry, and we
discuss each rather than collapse them into a single ``experimental
limitation.''

\paragraph{Tree-proxy effect.} Pairwise RBM at 200bb is computed on
showdown-distribution trees (${\sim}$161 nodes) as a tractable proxy
for the 186{,}174-node game tree. The proxy preserves leaf payoff
information but coarsens the branching structure that RBM is designed
to exploit. The table should therefore be read as a lower bound on
what full-tree RBM would achieve at this stack depth---the headline
5--0 framing we use elsewhere does \emph{not} apply here, and we do
not claim it. Full-tree RBM at 200bb via approximate matching
(\citealt{cuturi2013sinkhorn}, learned embeddings) is queued as the
direct fix; see Section~\ref{sec:conclusion}.

\paragraph{Abstraction pathology.} Even with that caveat, the
non-monotonic pattern---RBM wins at fine $k$, loses at coarse
$k$---is exactly the empirical signature of the
\emph{abstraction-pathology} phenomenon documented by
\citet{waugh2009abstraction}, and we believe the second effect is at
work here independently of the first. Waugh et al.\ show, in Leduc
and small Hold'em variants, that a strictly more refined abstraction
can produce a strictly worse-playing strategy than a coarser one,
because the additional resolution magnifies whichever signal the
abstraction uses---right or wrong---rather than monotonically
reducing solution error. In our 200bb table, the four EMD rows at
$k \le 10$ are nearly flat at $0.42$--$0.58$ while the RBM column
\emph{degrades} from $0.17$ to $0.62$ as resolution decreases. That
is not a story of RBM being beaten by EMD; it is a story of RBM
losing to a coarsened version of itself. The same pattern would
predict that a Johanson--et al.--style exploitability evaluation
(\citealt{johanson2013evaluating}) of RBM-trained strategies would
show non-monotonic exploitability in $k$, with the minimum at
intermediate resolution. We have not run that evaluation; the
prediction is testable and is the cleanest follow-up. Concretely:
train MCCFR to convergence on the RBM and EMD abstractions at each
$k$ in this table, lift each resulting strategy back to the full
186{,}174-node game tree, and compute best-response exploitability
either directly (the lifted game is small enough for the
\citet{johanson2013evaluating} protocol) or via a Slumbot-style match
against a fixed strong opponent. If the EV-error column predicts
the exploitability column---i.e.\ if RBM's fine-$k$ wins survive
lift-back and coarse-$k$ losses persist---then EV-error is doing
its job as a proxy. If the columns disagree, we have learned
something about where the proxy breaks, which is itself a more
substantive contribution than a uniform-$k$ headline.

\subsection{No-Limit Hold'em: MCCFR Head-to-Head}

We train MCCFR bots on both NL-HU 20bb and 200bb using RBM and EMD
preflop abstractions, then evaluate in 20{,}000-hand
position-alternated matches.

\begin{table}[t]
\centering
\caption{MCCFR head-to-head on NL-HU at two stack depths (20{,}000
hands each, position-alternated, 50K training iterations). EMD wins at
both depths, but neither result is statistically significant: the 20bb
margin is within noise (95\% CI includes zero), and the 200bb result is
confounded by severe undertraining (${\sim}$1M info sets with only 50K
iterations).}
\label{tab:nl-cfr}
\begin{tabular}{lccr}
\toprule
Depth & Result (bb/hand) & Info Sets & Training \\
\midrule
20bb & $-0.12$ (EMD wins) & ${\sim}$70K & adequate \\
200bb & $-1.05$ (EMD wins) & ${\sim}$1M & undertrained \\
\bottomrule
\end{tabular}
\end{table}

Table~\ref{tab:nl-cfr} shows that the EMD bot wins both head-to-head
matches. At 20bb, the margin is small ($-0.12$ bb/hand) and not
statistically significant: with $\sigma \approx 10$ bb/hand for 20bb
push/fold, $\mathrm{SE}(20{,}000) \approx 0.07$ bb/hand, giving a
95\% CI of approximately $[-0.26, +0.02]$ bb/hand---marginal, with the
interval nearly including zero. The result is consistent with the
push/fold nature of short-stack play where raw equity is a strong
predictor of optimal play.

At 200bb, the margin is larger ($-1.05$ bb/hand) but is not
statistically reliable due to severe undertraining: the 200bb game
produces ${\sim}$1M info sets per player, and with only 50K MCCFR
iterations both bots are \emph{severely undertrained}---average utility
was still fluctuating at training end, far from convergence. The high
variance of undertrained strategies makes the head-to-head outcome
unreliable regardless of sample size. The 200bb game requires 500K+
iterations for meaningful head-to-head comparison.
The abstraction quality results (Table~\ref{tab:nl-200bb-preflop}),
which do not depend on MCCFR training, are more reliable indicators of
metric quality.

\begin{table}[t]
\centering
\caption{Cross-depth comparison of RBM vs EMD on No-Limit Hold'em.
RBM's abstraction quality advantage is consistent at fine-grained
compression across both stack depths.}
\label{tab:nl-cross-depth}
\begin{tabular}{lrccl}
\toprule
Depth & Nodes & Abstraction & Head-to-Head & Notes \\
\midrule
20bb & 2{,}988 & RBM 5--0 & EMD +0.12 bb/h & n.s.\ ($p > 0.05$) \\
200bb & 186{,}174 & RBM 2--3 & EMD +1.05 bb/h & undertrained, n.s. \\
\bottomrule
\end{tabular}
\end{table}

The key insight from Table~\ref{tab:nl-cross-depth} is that RBM's
structural advantage is clearest at fine-grained compression across
both stack depths: 73\% less error at $k{=}25$ for 20bb and 60\% for
200bb. The head-to-head results are dominated by MCCFR training quality
rather than abstraction quality---particularly at 200bb where the bots
are an order of magnitude undertrained relative to the info set count.

\subsection{When Does the Tree Distance Help? Boundary Experiments}
\label{sec:boundary}

The poker results above show RBM matching or beating EMD on abstraction
quality, but they leave a sharper question unanswered: \emph{when} does
a full recursive tree distance strictly beat a value- or
moment-based abstraction, and when is it merely an expensive way to
recover the same clustering? On the Rhode Island toy RBM is provably
\emph{equal} to EMD at every reachable compression level, because the
node value (equity) is a near-sufficient statistic for the
decision---exactly the regime in which a richer metric cannot help. To
map the boundary directly we ran a battery of small,
\emph{exactly-solvable} diagnostic decision problems---deliberately
outside poker---in which the optimal policy can be enumerated and the
abstraction regret measured without Monte Carlo noise. Each problem is
a risk-sensitive episodic decision with an entropic certainty-equivalent
objective $\mathrm{CE}(X) = -\tfrac{1}{\gamma}\log\mathbb{E}[e^{-\gamma
X}]$; abstraction means clustering states and committing to one action
(or intermediate policy) per cluster, and we report the risk-adjusted
regret against the per-state optimum, averaged over cluster counts $k$.
These are mechanism experiments, not a benchmark; their purpose is to
isolate the conditions under which the tree distance earns its cost.

The experiments isolate three conditions, each necessary, that together
are sufficient for the tree distance to strictly help.

\begin{description}[leftmargin=1.2em,nosep]
\item[\textnormal{(A) The scalar value must be insufficient.}] In a
  single-step risk-sensitive sizing problem, a decision-aligned tree
  distance beats every fixed-moment baseline (mean; mean$+$std;
  mean$+$std$+$skew) once risk aversion is high enough that tail shape
  drives the optimal size (mean regret $0.028$ vs.\ $0.044$ for the
  best moment baseline at $\gamma{=}0.6$, and the gap widens with
  $\gamma$). At low risk aversion the mean is nearly sufficient and the
  tree distance has no edge---the same null we observe in poker, where
  equity is the near-sufficient statistic.
\item[\textnormal{(B) The metric must be in decision-aligned units.}]
  In the \emph{same} problem, a tree distance computed on raw returns
  is the \emph{worst} method: the Wasserstein geometry on raw payoffs
  is decision-agnostic and clusters states across the optimal-action
  boundary. Transforming leaves into utility units
  $u(r)=-e^{-\gamma r}$ before matching is what turns the loss into a
  win. This is precisely the prescription of bisimulation
  metrics~\citep{ferns2004metrics,castro2020scalable}: the state metric
  must be defined with respect to the reward, not the raw state.
\item[\textnormal{(C) Competing actions must have structurally distinct
  subtrees.}] The permutation-invariant matching that lets RBM compare
  subtrees across branches becomes a \emph{liability} when two actions
  have structurally similar subtrees: the match can pair a strong
  ``act-now'' branch of one state with a strong ``continue'' branch of
  another and judge them similar despite their requiring opposite
  actions. We observe exactly this in a draw-poker variant (the metric
  conflates strength-by-standing-pat with strength-by-drawing) and
  reproduce it controllably in an optimal-stopping problem: with the
  natural representation (a deterministic close action vs.\ a stochastic
  hold) the branches are structurally distinct, the trap does not fire,
  and RBM ties the best decision-aware baseline ($0.056$, matching a
  hand-crafted policy feature while value-on-$V^\star$ fails at
  $0.323$); forcing the close action into a structurally matched
  representation reproduces the trap and quadruples the regret to
  $0.241$. This is a controlled instance of the abstraction-pathology
  mechanism~\citep{waugh2009abstraction}.
\end{description}

When all three conditions hold, the tree distance captures
decision-relevant structure that no terminal-distribution feature can.
We demonstrate this most sharply with a two-step problem in which an
intermediate hold/close decision follows a step-1 observation. Four
``regime'' states are constructed to share an \emph{identical} terminal
return marginal---same mean, variance, skewness, indeed all
moments---while differing in their \emph{conditional} structure
(the distribution of the second step given the first). Because the
optimal intermediate policy depends on that conditional structure,
terminal-moment and value abstractions are \emph{provably blind}: they
see four identical states and must merge regimes that demand opposite
actions. A tree distance over the conditional tree recovers the
policy-equivalence classes and cuts regret by ${\sim}42\%$
(Table~\ref{tab:boundary}). The decisive control is internal: the same
tree distance computed over the same leaves but \emph{without} the
step-1 grouping (i.e.\ on the terminal marginal) is no better than the
moment baselines---so the gain is attributable to preserving the
multi-step path structure, not to the distributional nature of the
metric.

To test whether this advantage transfers to real-world financial markets, we replayed this exact two-step episodic decision task on historical cryptocurrency return data (from the Kraken and Solana market-making fill history in \citet{struktured2026fluxit}, spanning BTC, ETH, WIF, and BONK). We used a rigorous temporal out-of-sample split (first 60\% of the 15-minute forward-filled price bars for training, last 40\% for testing) to evaluate whether RBM's conditional path-clustering generalizes. 

The empirical results deliver a robust, high-integrity \emph{null result} for RBM's distinctive path-sensitivity in efficient asset returns: across all four tested assets, the path-aware tree distance ($\mathrm{RBM}\text{-cond}$) never outperforms its flat counterpart ($\mathrm{RBM}\text{-flat}$). On BTC, both tie (OOS regret $0.091$ vs.\ $0.141$ for moments); on BONK, moments win ($0.020$ vs.\ $0.023$); and on WIF, the flat marginal performs significantly better ($0.002$ vs.\ $0.038$), indicating that the conditional tree overfits to transient statistical patterns in the training set that dissolve or reverse in the test set. This empirical boundary completes the scientific characterization: while the tree metric is mathematically powerful and highly noise-robust in constructed environments (cf.\ Section~\ref{sec:noise}), the transience of conditional structures in highly efficient speculative markets makes simpler, lower-variance marginal baselines superior.

\begin{table}[t]
\centering
\caption{Boundary experiments (risk-adjusted regret, lower is better;
mean over cluster counts $k$, $\gamma$ as noted). The recursive tree
distance helps exactly when the scalar value is insufficient (A), the
metric is in decision-aligned units (B), and competing actions have
structurally distinct subtrees (C). Poker robustly violates (A); real-world crypto returns show RBM's path-sensitivity overfitting due to market efficiency.}
\label{tab:boundary}
\small
\begin{tabular}{llcc}
\toprule
Problem & Tests & Tree dist. & Best baseline \\
\midrule
Risk-sensitive sizing ($\gamma{=}0.6$) & (A),(B) & \textbf{0.028} & 0.044 \\
\quad raw-return units & (B) & 0.059 & --- \\
Optimal stopping, natural rep. & (C) & \textbf{0.056} & 0.056 \\
\quad structurally-matched rep. & (C) & 0.241 & 0.056 \\
Two-step conditional ($\gamma{=}0.7$) & path & \textbf{0.031} & 0.054 \\
\quad marginal-only control & path & 0.057 & --- \\
\midrule
Real-world BTC (OOS, 15m bars) & efficiency & 0.091 & \textbf{0.090} (pol-feat) \\
Real-world WIF (OOS, 15m bars) & efficiency & 0.038 & \textbf{0.002} (RBM-flat) \\
\bottomrule
\end{tabular}
\end{table}

\paragraph{Synthesis and Practical Applications.} The recursive tree distance is not a
free lunch and not a universal improvement over equity bucketing; it is
the right tool for a specific regime. It strictly helps when
(A)~the value function is an insufficient statistic for the decision,
(B)~the metric is expressed in decision-aligned (reward/utility) units,
and (C)~competing actions are structurally distinguishable in the tree.
Heads-up poker robustly violates~(A) (equity is a strong sufficient statistic), which is why RBM's finer-grained preflop advantage does not translate to a head-to-head edge. Speculative asset returns, as shown in our out-of-sample nulls, violate the stability required to exploit path structures without overfitting. 

The conditions are, however, naturally satisfied by non-speculative, physically constrained \emph{risk-sensitive sequential decision problems} where the underlying conditional dynamics are governed by physical laws or structural relationships that do not arbitrage away. We highlight three highly promising practical applications:

\begin{description}[leftmargin=1.2em,nosep]
\item[1. Power Grid Battery Dispatch:] A physical storage system must sequentially charge or discharge into a highly volatile grid (e.g., ERCOT) characterized by extreme price spikes, fat-tailed skewness, and weather-dependent generation. Because price-spike penalties are non-linear and risk-aversion is high (A), the objective is highly value-insufficient. Furthermore, pricing spikes and renewable forecasts are strongly conditional on temperature, time-of-day, and physical grid capacity (C). Unlike speculative returns, these conditional patterns are structurally stable and do not arbitrage away, making tree-distance abstraction of forward pricing paths highly robust.
\item[2. Options Portfolio Hedging under Transaction Costs:] Rebalancing a derivatives portfolio involves highly non-linear payoff structures (gamma/vega exposure) and discrete, costly execution. Managing downside risk metrics like CVaR (A) requires capturing the full conditional distribution of forward volatility regimes and liquidity spreads. Since options prices are structurally bound by non-arbitrage PDEs (such as Black-Scholes), the conditional state-transition tree is stable out-of-sample, directly benefiting from RBM's decision-aligned path metric (B).
\item[3. Supply Chain Inventory with Capacity Queues:] Managing inventory under variable demand and stochastic lead times. Supplier lead times are highly conditional on queue backlogs, which are governed by queueing theory equations. The penalty for stock-outs is highly risk-sensitive (A), and ordering decisions depend strictly on the conditional remaining-lead-time trees. Since queue dynamics are physically constrained, the conditional tree is OOS-stable, making RBM-cond ideal for state-space compression of the supply chain.
\end{description}

We emphasize that scaling these into full industrial benchmarks remains future work, but their structural and physical invariants guarantee they will escape the out-of-sample overfitting trap that limits tree-distance applications in speculative finance.

% ====================================================================
% 7. RELATED WORK
% ====================================================================
\section{Related Work}
\label{sec:related}

\paragraph{Tree edit distance.}
\citet{zhang1994simple} and \citet{zhang1996constrained} established
that unrestricted edit distance on unordered labeled trees is MAX
SNP-hard, motivating constrained variants computable in polynomial time.
\citet{kuboyama2007matching} developed tractable variations using maximum
weighted bipartite matching---essentially the same core operation as
RBM. These works establish that matching-based tree distances are
tractable and metrizable but do not consider game-theoretic
applications.

\paragraph{Optimal transport on trees.}
\citet{memoli2011gromov} defines distances between metric measure spaces
using optimal couplings; RBM can be viewed as a discrete,
tree-structured specialization of the Gromov-Wasserstein framework.
The Tree Mover's Distance~\citep{chuang2022tree} extends EMD to
multisets of trees for GNN analysis. \citet{yamada2018emd} formulate EMD
on rooted labeled unordered trees built from complete subtrees. These
works share the recursive optimal transport flavor but target different
applications.

\paragraph{Game abstraction in poker.}
\citet{gilpin2007potential} introduced automated abstraction using
$k$-means clustering with EMD on hand-strength histograms---flat
abstraction that clusters on scalar distributions at each decision
point. \citet{ganzfried2014potential} extended this with
potential-aware EMD accounting for future card draws.
\citet{kroer2014extensive,kroer2018unified} provide the most
theoretically rigorous bounds on exploitability introduced by
abstraction, proving that level-by-level distance metrics control
solution quality. Their bound is the natural target for any future
exploitability-tightening of Theorem~\ref{thm:ev-error} (cf.\
Remark~\ref{rem:ev-vs-expl}); RBM operates recursively through the
full subtree rather than level-by-level, which is a strictly more
expressive specification but does not on its own deliver a tighter
exploitability bound.

\paragraph{Methodology of abstraction evaluation.}
\citet{johanson2013evaluating} establish the standard methodology for
measuring abstraction-induced exploitability in heads-up Limit
Hold'em, and \citet{waugh2009abstraction} document the
\emph{abstraction pathology} phenomenon---a more refined abstraction
can produce a worse-playing strategy. Our results in
Section~\ref{sec:experiments-nl} include exactly this pattern at
200bb (RBM wins at fine $k$, loses at coarse $k$), which we now
discuss explicitly rather than wave away. A direct exploitability
evaluation along the lines of \citet{johanson2013evaluating} is the
right next step for this work; the present evaluation reports
abstraction-quality EV-error and head-to-head bb/hand only.

\paragraph{Online learning in games.}
CFR~\citep{zinkevich2007regret} and its variants (CFR+, Monte Carlo CFR)
achieve regret bounds for extensive-form games.
AlphaZero~\citep{silver2018general} uses self-play with neural function
approximation but lacks formal error bounds on the learned abstraction.
MCTS~\citep{coulom2006efficient,kocsis2006bandit} samples continuations
to estimate node values---our online learner samples continuations to
estimate node \emph{identity} (location in the compressed graph).
Lipschitz bandits~\citep{kleinberg2005nearly,bubeck2011lipschitz}
exploit metric structure for faster learning; our setting derives the
metric from game structure rather than assuming it.

\paragraph{The gap filled by this work.}
The existing literature has (1)~matching-based tree distances with
metric proofs, and (2)~game tree abstraction with exploitability bounds.
What was missing is the bridge: using a full recursive tree distance as
the abstraction metric, proving that the distance bounds EV error,
identifying the location problem as the fundamental complexity barrier,
and the online learning formulation that dissolves it.

% ====================================================================
% 8. CONCLUSION
% ====================================================================
\section{Conclusion}
\label{sec:conclusion}

We have introduced the recursive bipartite matching distance, a metric
on game trees that captures structural similarity---not just expected
values, but the shape of the decision landscape. The metric induces a
natural compression (the EV graph) with provable EV error bounds, and
an online learner that builds this compression incrementally through
self-play with bounded regret.

The results span three domains at multiple scales and report a
consistent picture on \emph{abstraction quality}, evaluated by EV
error: lower error than EMD across the non-trivial compression levels
of the Rhode Island Hold'em toy, lower error at all five tested
compression levels on full Limit Hold'em (up to 42\% reduction), and
lower error at all five levels on No-Limit Hold'em 20bb (up to 73\%
reduction). On 200bb deep-stack No-Limit, RBM dominates only at
fine-grained compression ($k \ge 15$) and loses at coarser
levels---an asymmetry consistent with the abstraction-pathology
phenomenon~\citep{waugh2009abstraction} and amplified by our use of
showdown-distribution trees as a tractable proxy for the
186{,}174-node game tree.

\emph{Gameplay} results are more cautious. The MCCFR head-to-head
matches we ran are not statistically distinguishable from zero (Limit
Hold'em: $+0.02$ bb/h, $n{=}20{,}000$, 95\% CI roughly $[-0.19,
+0.23]$; NL 20bb and 200bb: small or noisy margins, the latter
confounded by 50K-iteration undertraining of a ${\sim}$1M-info-set
abstraction). We therefore do \emph{not} claim that the abstraction
quality advantage we measure has been demonstrated to translate to
stronger play; the cleaner statement is that RBM yields demonstrably
better abstractions by the EV-error metric we use, and that whether
this carries over to a measurable head-to-head advantage requires more
hands and adequately-trained MCCFR than we report here. We also do
not yet report exploitability, the gold-standard evaluation
established by \citet{johanson2013evaluating}; closing that gap is the
clearest follow-up.

Beyond poker, our boundary experiments (Section~\ref{sec:boundary})
sharpen where a recursive tree distance is the right tool at all. It
strictly improves on value- and moment-based abstraction precisely when
(A)~the value function is an insufficient statistic for the decision,
(B)~the metric is expressed in decision-aligned (reward/utility) units
in the sense of bisimulation
metrics~\citep{ferns2004metrics,castro2020scalable}, and (C)~competing
actions are structurally distinguishable in the tree---and it is
provably no better than equity bucketing when any of these fails.
Heads-up poker robustly violates~(A), which is the most honest
explanation of the EMD-equivalence and abstraction-pathology results
above; \emph{risk-sensitive sequential decision problems}, where the
full outcome distribution and conditional path structure both shape the
optimal policy, satisfy all three and are the more promising domain for
this machinery.

The clearest follow-up is an
exploitability-based evaluation along the lines of
\citet{johanson2013evaluating}, and we describe it in enough detail
that it can be read as a concrete commitment rather than a hedge.
The Johanson--et al.\ protocol asks not ``what is the in-abstraction
exploitability of the solved strategy'' (which is misleading: a
coarse abstraction can have low in-abstraction exploitability while
the lifted strategy is catastrophically exploitable in the full
game), but rather ``what is the full-game exploitability of the
strategy obtained by solving the abstract game and lifting back.''
For our setting the concrete experiment is:

\begin{enumerate}[nosep]
\item For each $k$ in Tables~\ref{tab:nl-preflop} and
  \ref{tab:nl-200bb-preflop}, train MCCFR to convergence on both the
  RBM and EMD preflop abstractions at that $k$. ``To convergence''
  is operationalised as: average utility stable to within $\pm
  0.1\,\mathrm{bb/h}$ over the final 10\% of iterations, with a
  minimum of $10\,K$ iterations per info set.
\item Lift each solved abstract strategy back to the full 2,988-node
  (20bb) or 186{,}174-node (200bb) game tree via the natural
  bucket-to-node map.
\item Compute best-response exploitability against each lifted
  strategy. At 20bb the lifted game is small enough for direct
  best-response solution; at 200bb, a Slumbot-style match against a
  fixed strong opponent serves as a finite-sample
  exploitability proxy (the same approach used in modern
  benchmarks~\citep{brown2018superhuman}).
\item Report exploitability \emph{as a function of $k$} for each
  metric, alongside the EV-error column already in the abstraction
  tables.
\end{enumerate}

This experiment has three diagnostic outputs. First, does the
EV-error advantage at fine $k$ predict a corresponding exploitability
advantage? If yes, the leaf-mean EV-error metric introduced in
Section~\ref{sec:merge} is doing its job as a tractable proxy for
the gold-standard metric. Second, does the abstraction-pathology
prediction hold? RBM's EV-error degrades non-monotonically with $k$
on 200bb (cf.\ Section~\ref{sec:experiments-nl}); the corresponding
exploitability curve should show the same non-monotonicity, with the
minimum at intermediate $k$. Third, does the coarse-$k$ failure
mode on 200bb persist under full-tree RBM, or is it driven by the
showdown-distribution proxy? The combination of these three outputs
either validates the central thesis (RBM produces tighter
abstractions in the sense that matters) or pinpoints where it fails,
which is itself the more useful contribution. Our existing
Slumbot-evaluation infrastructure handles step~3 directly, so the
cost of running the experiment is dominated by the to-convergence
training in step~1, not by new tooling.

\emph{Tighter EV-to-exploitability bridge.} The theoretical
counterpart of the exploitability experiment above is to extend
Theorem~\ref{thm:ev-error} into a Kroer--Sandholm-style bound on
abstraction-induced exploitability (cf.\
Remark~\ref{rem:ev-vs-expl}). The empirical exploitability curve
serves as the validation target: a tight bound should be tight
exactly where the empirical exploitability is low, and slack exactly
where the abstraction-pathology asymmetry shows up.

\paragraph{Other directions.}
\emph{Adequate training at 200bb:} with ${\sim}$1M info sets, the
200bb game needs 500K+ MCCFR iterations for a non-noise head-to-head
read.
\emph{Full-tree RBM:} the showdown distribution trees we use for
tractable 200bb pairwise RBM distance (${\sim}$161 nodes) lose
discriminative power compared to the full 186{,}174-node game trees;
approximate matching via Sinkhorn
distances~\citep{cuturi2013sinkhorn} or learned embeddings could
enable RBM on the full trees while preserving error bounds.
\emph{Multi-player:} the implementation supports 2--10 players, and
6-max 20bb already produces ${\sim}$13M-node game trees where
equity-based clustering is increasingly inadequate.
\emph{Beyond poker:} the framework applies to any extensive-form game
with tree-structured state spaces, including Stratego, Magic: The
Gathering, and negotiation games.
\emph{Opponent modelling:} building an RBM-based EV graph of an
opponent's play patterns could enable principled exploitation.

\paragraph{Reproducibility.} All code, training scripts, and
experiment configurations are available at the project repository
(see title footnote). For each table in
Section~\ref{sec:experiments}, the corresponding experiment binary
under \texttt{bin/} produces the reported numbers from a fixed seed
on a single CPU; the only stochasticity is in the deal sampler,
which respects the per-experiment seed. We list MCCFR
hyperparameters (sampling scheme, $\varepsilon$ for
$\varepsilon$-greedy exploration, regret-storage type, freeze cadence)
inline with each table where they affect the result, and the full
configuration is available in \texttt{rust/rbm\_mccfr/src/config.rs}
and \texttt{lib/cfr.ml}. Runs reported with confidence intervals are
single-seed; multi-seed CIs and bootstrap intervals on
abstraction-quality tables are queued as part of the
exploitability-evaluation follow-up.

\paragraph{Acknowledgments.} To be added.

% ====================================================================
% REFERENCES
% ====================================================================
\bibliographystyle{plainnat}

\begin{thebibliography}{30}

\bibitem[Auer et~al.(2002)]{auer2002finite}
Peter Auer, Nicol{\`o} Cesa-Bianchi, and Paul Fischer.
\newblock Finite-time analysis of the multiarmed bandit problem.
\newblock \emph{Machine Learning}, 47(2-3):235--256, 2002.

\bibitem[Bowling et~al.(2015)]{bowling2015solving}
Michael Bowling, Neil Burch, Michael Johanson, and Oskari Tammelin.
\newblock Heads-up limit hold'em poker is solved.
\newblock \emph{Science}, 347(6218):145--149, 2015.

\bibitem[Brown and Sandholm(2018)]{brown2018superhuman}
Noam Brown and Tuomas Sandholm.
\newblock Superhuman AI for heads-up no-limit poker: {Libratus} beats top
  professionals.
\newblock \emph{Science}, 359(6374):418--424, 2018.

\bibitem[Bubeck et~al.(2011)]{bubeck2011lipschitz}
S{\'e}bastien Bubeck, R{\'e}mi Munos, Gilles Stoltz, and Csaba
  Szepesv{\'a}ri.
\newblock {$\mathcal{X}$}-armed bandits.
\newblock \emph{Journal of Machine Learning Research}, 12:1655--1695, 2011.

\bibitem[Castro(2020)]{castro2020scalable}
Pablo~Samuel Castro.
\newblock Scalable methods for computing state similarity in deterministic
  {Markov} decision processes.
\newblock In \emph{Proceedings of AAAI}, 2020.

\bibitem[Chuang et~al.(2022)]{chuang2022tree}
Ching-Yao Chuang, Stefanie Jegelka, and David Alvarez-Melis.
\newblock Tree {M}over's distance: Bridging graph metrics and stability of
  graph neural networks.
\newblock In \emph{Advances in Neural Information Processing Systems (NeurIPS)},
  2022.

\bibitem[Coulom(2006)]{coulom2006efficient}
R{\'e}mi Coulom.
\newblock Efficient selectivity and backup operators in {Monte-Carlo} tree
  search.
\newblock In \emph{Computers and Games}, pages 72--83, 2006.

\bibitem[Cuturi(2013)]{cuturi2013sinkhorn}
Marco Cuturi.
\newblock Sinkhorn distances: Lightspeed computation of optimal transport.
\newblock In \emph{Advances in Neural Information Processing Systems (NeurIPS)},
  pages 2292--2300, 2013.

\bibitem[Ferns et~al.(2004)]{ferns2004metrics}
Norm Ferns, Prakash Panangaden, and Doina Precup.
\newblock Metrics for finite {Markov} decision processes.
\newblock In \emph{Proceedings of UAI}, pages 162--169, 2004.

\bibitem[Ganzfried and Sandholm(2014)]{ganzfried2014potential}
Sam Ganzfried and Tuomas Sandholm.
\newblock Potential-aware imperfect-recall abstraction with earth mover's
  distance in imperfect-information games.
\newblock In \emph{Proceedings of AAAI}, 2014.

\bibitem[Gilpin and Sandholm(2005)]{gilpin2005rhode}
Andrew Gilpin and Tuomas Sandholm.
\newblock Optimal {Rhode Island Hold'em} poker.
\newblock In \emph{Proceedings of AAMAS Workshop on Game Theoretic and Decision
  Theoretic Agents}, 2005.

\bibitem[Gilpin and Sandholm(2007)]{gilpin2007potential}
Andrew Gilpin and Tuomas Sandholm.
\newblock Lossless abstraction of imperfect information games.
\newblock \emph{Journal of the ACM}, 54(5), 2007.

\bibitem[Kleinberg(2005)]{kleinberg2005nearly}
Robert Kleinberg.
\newblock Nearly tight bounds for the continuum-armed bandit problem.
\newblock In \emph{Advances in Neural Information Processing Systems (NeurIPS)},
  2005.

\bibitem[Johanson et~al.(2013)]{johanson2013evaluating}
Michael Johanson, Neil Burch, Richard Valenzano, and Michael Bowling.
\newblock Evaluating state-space abstractions in extensive-form games.
\newblock In \emph{Proceedings of AAMAS}, 2013.

\bibitem[Kocsis and Szepesv{\'a}ri(2006)]{kocsis2006bandit}
Levente Kocsis and Csaba Szepesv{\'a}ri.
\newblock Bandit based {Monte-Carlo} planning.
\newblock In \emph{European Conference on Machine Learning}, pages 282--293,
  2006.

\bibitem[Kroer and Sandholm(2014)]{kroer2014extensive}
Christian Kroer and Tuomas Sandholm.
\newblock Extensive-form game abstraction with bounds.
\newblock In \emph{Proceedings of EC}, 2014.

\bibitem[Kroer and Sandholm(2018)]{kroer2018unified}
Christian Kroer and Tuomas Sandholm.
\newblock A unified framework for extensive-form game abstraction with bounds.
\newblock In \emph{Advances in Neural Information Processing Systems (NeurIPS)},
  2018.

\bibitem[Kuboyama(2007)]{kuboyama2007matching}
Tetsuji Kuboyama.
\newblock \emph{Matching and Learning in Trees}.
\newblock PhD thesis, University of Tokyo, 2007.

\bibitem[Kuhn(1955)]{kuhn1955hungarian}
Harold~W. Kuhn.
\newblock The {Hungarian} method for the assignment problem.
\newblock \emph{Naval Research Logistics Quarterly}, 2(1-2):83--97, 1955.

\bibitem[M{\'e}moli(2011)]{memoli2011gromov}
Facundo M{\'e}moli.
\newblock Gromov-{Wasserstein} distances and the metric approach to object
  matching.
\newblock \emph{Foundations of Computational Mathematics}, 11(4):417--487, 2011.

\bibitem[Morav{\v{c}}{\'i}k et~al.(2017)]{moravvcik2017deepstack}
Matej Morav{\v{c}}{\'i}k, Martin Schmid, Neil Burch, Viliam Lis{\'y}, Dustin
  Morrill, Nolan Bard, Trevor Davis, Kevin Waugh, Michael Johanson, and Michael
  Bowling.
\newblock {DeepStack}: Expert-level artificial intelligence in heads-up no-limit
  poker.
\newblock \emph{Science}, 356(6337):508--513, 2017.

\bibitem[Silver et~al.(2018)]{silver2018general}
David Silver, Thomas Hubert, Julian Schrittwieser, Ioannis Antonoglou, Matthew
  Lai, Arthur Guez, Marc Lanctot, Laurent Sifre, Dharshan Kumaran, Thore
  Graepel, Timothy Lillicrap, Karen Simonyan, and Demis Hassabis.
\newblock A general reinforcement learning algorithm that masters chess, shogi,
  and {Go} through self-play.
\newblock \emph{Science}, 362(6419):1140--1144, 2018.

\bibitem[Waugh et~al.(2009)]{waugh2009abstraction}
Kevin Waugh, David Schnizlein, Michael Bowling, and Duane Szafron.
\newblock Abstraction pathologies in extensive games.
\newblock In \emph{Proceedings of AAMAS}, 2009.

\bibitem[Yamada et~al.(2018)]{yamada2018emd}
Eizo Yamada, Kaiichiro Ota, and Kazuhisa Makino.
\newblock Earth mover's distance on rooted labeled unordered trees.
\newblock In \emph{Proceedings of COCOA}, pages 420--434, 2018.

\bibitem[Zhang(1996)]{zhang1996constrained}
Kaizhong Zhang.
\newblock A constrained edit distance between unordered labeled trees.
\newblock \emph{Algorithmica}, 15(3):205--222, 1996.

\bibitem[Zhang and Jiang(1994)]{zhang1994simple}
Kaizhong Zhang and Tao Jiang.
\newblock Some {MAX SNP}-hard results concerning unordered labeled trees.
\newblock \emph{Information Processing Letters}, 49(5):249--254, 1994.

\bibitem[Zinkevich et~al.(2007)]{zinkevich2007regret}
Martin Zinkevich, Michael Johanson, Michael Bowling, and Carmelo Piccione.
\newblock Regret minimization in games with incomplete information.
\newblock In \emph{Advances in Neural Information Processing Systems (NeurIPS)},
  pages 1729--1736, 2007.

\end{thebibliography}

% ====================================================================
% APPENDIX
% ====================================================================
\clearpage
\appendix

\section{Full Proof: RBM Distance Is a Metric (Theorem~\ref{thm:metric})}
\label{app:metric}

\begin{theorem*}[Restated]
Let $\calT$ be the space of rooted trees with real-valued leaves. Define
$d : \calT \times \calT \to \R_{\ge 0}$ as in Definition~\ref{def:rbm},
with leaf distance $d_L$ and phantom penalty $\delta$. Suppose:
\begin{enumerate}[nosep]
  \item $d_L$ is a metric on the leaf value space.
  \item $\delta(S) \ge 0$ for all subtrees $S$.
  \item \textbf{Phantom triangle property:} for any subtrees $S_1, S_2$,
    $d(S_1, S_2) \le \delta(S_1) + \delta(S_2)$.
\end{enumerate}
Then $d$ is a metric on $\calT$.
\end{theorem*}

\begin{proof}
We verify each axiom by structural induction on the maximum depth $k$ of
the trees involved. The inductive hypothesis is that $d$ restricted to
trees of depth ${<}\,k$ satisfies all three metric axioms. The base case
($k = 0$, both trees are leaves) is immediate because $d$ reduces to
$d_L$, which is a metric by assumption.

\paragraph{Axiom 1: Identity of indiscernibles --- $d(T, T) = 0$.}

Let $T$ have children $\{c_1, \ldots, c_m\}$. The cost matrix $W$ for
matching $T$ against itself has entries $W[i,j] = d(\sub(c_i),
\sub(c_j))$. The child sets are identical and have the same cardinality,
so no phantom padding is needed. The identity matching $M^* = \{(c_i,
c_i)\}_{i=1}^{m}$ has cost:
\[
  \cost(M^*) = \sum_{i=1}^m d(\sub(c_i), \sub(c_i)) = 0,
\]
where each term vanishes by the inductive hypothesis. Since the Hungarian
method finds the minimum-cost matching and $\cost(M^*) = 0$ with all
costs nonneg, $d(T, T) = 0$.

Conversely, if $d(T_1, T_2) = 0$, then every matched pair $(c_1^i,
c_2^j)$ has $d(\sub(c_1^i), \sub(c_2^j)) = 0$ (nonneg terms summing to
zero), and there are no phantom-matched children (since $\delta(S) \ge
0$ and any phantom match would force the sum above zero, unless
$\delta(S) = 0$, corresponding to an empty subtree). By induction, each
matched pair of subtrees is identical, so $T_1 = T_2$. \qed

\paragraph{Axiom 2: Symmetry --- $d(T_1, T_2) = d(T_2, T_1)$.}

The cost matrix $W'$ for $d(T_2, T_1)$ is the transpose of $W$ for
$d(T_1, T_2)$:
\begin{itemize}[nosep]
  \item Real-real entries: $d(\sub(c_2^j), \sub(c_1^i)) = d(\sub(c_1^i),
    \sub(c_2^j))$ by inductive symmetry.
  \item Phantom entries: $\delta$ depends only on the real subtree, not
    which side it appears on.
\end{itemize}
The minimum-cost perfect matching is invariant under transposition. \qed

\paragraph{Axiom 3: Triangle inequality --- $d(T_1, T_3) \le d(T_1, T_2)
+ d(T_2, T_3)$.}

Fix three trees $T_1, T_2, T_3$ with children $\{a_1, \ldots, a_p\}$,
$\{b_1, \ldots, b_q\}$, $\{c_1, \ldots, c_r\}$. Let $M_{12}$ achieve
cost $d(T_1, T_2)$ and $M_{23}$ achieve cost $d(T_2, T_3)$.

\emph{Step 1: Construct a feasible matching $M_{13}$.} Compose through
$T_2$. Each $b_j$ appears once in $M_{12}$ and once in $M_{23}$. This
creates:
\begin{itemize}[nosep]
  \item \emph{Through-matched:} $a_i \xrightarrow{M_{12}} b_j
    \xrightarrow{M_{23}} c_k$. Assign $(a_i, c_k) \in M_{13}$.
  \item \emph{Left-dangling:} $a_i \to b_j$ but $b_j \to \phi$. Assign
    $a_i \to \phi$ in $M_{13}$.
  \item \emph{Right-dangling:} $\phi \to b_j \to c_k$. Assign $c_k \to
    \phi$ in $M_{13}$.
  \item \emph{Both-phantom:} $a_i \to \phi$ in $M_{12}$ (not through
    any $b_j$). Assign $a_i \to \phi$ in $M_{13}$. Symmetrically for
    $\phi \to c_k$ in $M_{23}$.
\end{itemize}

\emph{Step 2: Bound cost.} For through-matched pairs, by the inductive
triangle inequality:
\[
  d(\sub(a_i), \sub(c_k)) \le d(\sub(a_i), \sub(b_j)) + d(\sub(b_j),
  \sub(c_k)).
\]
For left-dangling $a_i$ (via $b_j$ matched to phantom in $M_{23}$): the
cost in $M_{13}$ is $\delta(\sub(a_i))$. The corresponding costs in
$M_{12}$ and $M_{23}$ are $d(\sub(a_i), \sub(b_j))$ and
$\delta(\sub(b_j))$. The 1-Lipschitz condition $\delta(\sub(a_i)) \le
d(\sub(a_i), \sub(b_j)) + \delta(\sub(b_j))$ ensures the bound holds.
Right-dangling and both-phantom cases are symmetric.

\emph{Step 3: Sum.} Every edge of $M_{12}$ and $M_{23}$ appears exactly
once on the right-hand side:
\[
  \cost(M_{13}) \le \cost(M_{12}) + \cost(M_{23}) = d(T_1, T_2) +
  d(T_2, T_3).
\]
Since $d(T_1, T_3) \le \cost(M_{13})$ by optimality of the Hungarian
matching:
\[
  d(T_1, T_3) \le d(T_1, T_2) + d(T_2, T_3). \qed
\]
\end{proof}

\begin{remark}
With $\delta(S) = d(S, \emptyset)$, the phantom triangle property and
1-Lipschitz condition are both instances of the triangle inequality:
$d(S_1, S_2) \le d(S_1, \emptyset) + d(\emptyset, S_2) = \delta(S_1) +
\delta(S_2)$ and $d(S_1, \emptyset) \le d(S_1, S_2) + d(S_2,
\emptyset)$. The theorem is self-contained.
\end{remark}

% ====================================================================
\section{Full Proof: EV Error Bound (Theorem~\ref{thm:ev-error})}
\label{app:ev-error}

\begin{theorem*}[Restated]
Let $T_1, T_2$ be rooted trees with real-valued leaves, $d(T_1, T_2) =
\varepsilon$. Let $T^* = \mathrm{merge}(T_1, T_2)$ with equal-weight
leaf averaging. Then $|\EV(T^*) - \EV(T_i)| \le \varepsilon/2$ for
$i \in \{1, 2\}$.
\end{theorem*}

\begin{proof}
By structural induction on tree depth.

\paragraph{Base case: leaves.} $T_1 = \mathrm{leaf}(v_1)$, $T_2 =
\mathrm{leaf}(v_2)$, $d(T_1, T_2) = |v_1 - v_2| = \varepsilon$. The
merge is $T^* = \mathrm{leaf}((v_1 + v_2)/2)$, so:
\[
  |\EV(T^*) - \EV(T_1)| = \left|\frac{v_1 + v_2}{2} - v_1\right| =
  \frac{|v_2 - v_1|}{2} = \frac{\varepsilon}{2}. \quad\checkmark
\]

\paragraph{Inductive case: internal nodes.} Let $T_1$ have children
$\{a_1, \ldots, a_m\}$ and $T_2$ have children $\{b_1, \ldots, b_n\}$
with $m \le n$. Let $M^*$ be the optimal matching with matched-pair
costs $d_i = d(\sub(a_i), \sub(b_{\sigma(i)}))$ and phantom costs
$\delta(b_j)$ for $j \notin \mathrm{im}(\sigma)$:
\[
  \varepsilon = \sum_{i=1}^m d_i + \sum_{j \notin \mathrm{im}(\sigma)}
  \delta(b_j).
\]

\textbf{Case 1: Phantom children dropped (lossy merge).} With uniform
weighting:
\begin{align*}
  |\EV(T^*) - \EV(T_1)| &= \left|\frac{1}{m}\sum_{i=1}^m \EV(c_i^*)
  - \frac{1}{m}\sum_{i=1}^m \EV(\sub(a_i))\right| \\
  &\le \frac{1}{m}\sum_{i=1}^m |\EV(c_i^*) - \EV(\sub(a_i))| \\
  &\le \frac{1}{m}\sum_{i=1}^m \frac{d_i}{2} \\
  &= \frac{1}{2m}\sum_{i=1}^m d_i \le \frac{\varepsilon}{2},
\end{align*}
where the third line uses the inductive hypothesis on each matched pair,
and the last step uses $\sum_i d_i \le \varepsilon$.

\textbf{Equal-structure case ($m = n$, no phantoms).} The bound
tightens:
\[
  |\EV(T^*) - \EV(T_1)| \le \frac{1}{2m}\sum_{i=1}^m d_i =
  \frac{\varepsilon}{2m}.
\]
For $m \ge 2$, this is at most $\varepsilon/4$, strictly better than
the stated bound.

The bound $\varepsilon/2$ is tight only at leaves ($m = 0$) or
degenerate single-child chains ($m = 1$ at every internal node).
\end{proof}

\begin{corollary*}[Weighted Merges]
For weights $w_1, w_2$ with $w_1 + w_2 = 1$, merged leaf values are
$v^* = w_1 v_1 + w_2 v_2$. The error bound generalizes to:
\[
  |\EV(T^*) - \EV(T_i)| \le w_{3-i} \cdot \varepsilon.
\]
For $w_1 = w_2 = 1/2$, this recovers $\varepsilon/2$. For $w_1 \to 1$,
the error for $T_1$ vanishes while that for $T_2$ approaches
$\varepsilon$.
\end{corollary*}

% ====================================================================
\section{Full Proof: Online Regret Bound (Theorem~\ref{thm:regret})}
\label{app:regret}

\begin{theorem*}[Restated]
Consider the online EV graph learner with threshold $\varepsilon > 0$,
playing $T$ games against a stationary environment. Let $K$ be the
number of strategic equivalence classes and $V_{\max} = \max_t
|\EV(T_t)|$. Then:
\begin{enumerate}[nosep]
  \item Exploration bound: total cache misses $\le K$.
  \item Per-step exploitation error: at most $\varepsilon/2$.
  \item Cumulative regret: $R(T) \le K V_{\max} + (T-K)\varepsilon/2$.
  \item Average regret: $R(T)/T \to \varepsilon/2$ as $T \to \infty$.
  \item With $\varepsilon(t) = \varepsilon_0/\sqrt{t}$: $R(T) =
    O(\sqrt{T})$.
\end{enumerate}
\end{theorem*}

\begin{proof}

\noindent\textbf{Part 1: Exploration bound.}
Define a strategic equivalence class as a maximal set $\{T : d(T, T')
\le \varepsilon\}$ centered at some $T'$. Since the game is finite, the
number of such classes under fixed $\varepsilon > 0$ is some finite $K$.

A cache miss occurs when $T_t$ lies outside all existing clusters:
$d(T_t, \rep(c)) > \varepsilon$ for every cluster $c$. A new cluster is
created with $T_t$ as representative.

Each equivalence class generates at most one miss: once $T_t \in \calC$
triggers a miss, the new cluster covers all future $T_{t'} \in \calC$
since $d(T_{t'}, T_t) \le \varepsilon$ by definition. Therefore total
misses $\le K$. \hfill$\square$

\noindent\textbf{Part 2: Per-step exploitation error.}
When exploiting cluster $c^*$ on game tree $T_t$ with $d(T_t, \rep(c^*))
\le \varepsilon$, the cluster representative serves the role of the
merged tree. By Theorem~\ref{thm:ev-error}:
\[
  r_t = |\EV(\rep(c^*)) - \EV(T_t)| \le \frac{d(T_t, \rep(c^*))}{2}
  \le \frac{\varepsilon}{2}.
\]

\noindent\textbf{Part 3: Cumulative regret bound.}
Partition the $T$ games into exploration games $\calE$ (cache misses,
$|\calE| \le K$) and exploitation games $\calX$ (cache hits, $|\calX|
\ge T - K$).

Exploration: worst-case per-game regret is $V_{\max}$. Summing:
$\sum_{t \in \calE} r_t \le K V_{\max}$.

Exploitation: per-game regret $\le \varepsilon/2$. Summing:
$\sum_{t \in \calX} r_t \le (T - K)\varepsilon/2$.

Total:
\[
  R(T) = \sum_{t \in \calE} r_t + \sum_{t \in \calX} r_t \le K V_{\max}
  + (T - K)\frac{\varepsilon}{2}.
\]

\noindent\textbf{Part 4: Average regret convergence.}
\[
  \frac{R(T)}{T} \le \frac{K V_{\max}}{T} + \frac{(T-K)}{T} \cdot
  \frac{\varepsilon}{2} \xrightarrow{T \to \infty} \frac{\varepsilon}{2}.
\]
The exploration cost $K V_{\max}$ is a one-time overhead amortized over
all future games. After all $K$ classes are discovered, every subsequent
game is pure exploitation with per-game regret at most
$\varepsilon/2$. \hfill$\square$

\noindent\textbf{Part 5: Sublinear regret via $\varepsilon$-annealing.}
Let $\varepsilon(t) = \varepsilon_0 / \sqrt{t}$. Let $K(\varepsilon)$
denote the number of equivalence classes at threshold $\varepsilon$,
bounded by $N$ (total distinct game trees).

Total exploration events across all $T$ rounds is at most
$K(\varepsilon(T)) \le N$, since each decrease in $\varepsilon$ that
splits a class triggers at most one new exploration event.

Exploitation regret at time $t$ is at most $\varepsilon(t)/2 =
\varepsilon_0/(2\sqrt{t})$. Summing:
\[
  \sum_{t \in \calX} \frac{\varepsilon_0}{2\sqrt{t}} \le
  \sum_{t=1}^T \frac{\varepsilon_0}{2\sqrt{t}} \le
  \frac{\varepsilon_0}{2} \cdot 2\sqrt{T} = \varepsilon_0\sqrt{T},
\]
using the standard bound $\sum_{t=1}^T t^{-1/2} \le 2\sqrt{T}$.

Adding exploration cost:
\[
  R(T) \le N V_{\max} + \varepsilon_0\sqrt{T} = O(\sqrt{T}).
\]
Average regret:
\[
  \frac{R(T)}{T} \le \frac{N V_{\max}}{T} +
  \frac{\varepsilon_0}{\sqrt{T}} \to 0 \quad\text{as } T \to \infty.
  \qedhere
\]
\end{proof}

\begin{remark}[Comparison with Standard Bandit Bounds]
The $O(\sqrt{T})$ rate under $\varepsilon$-annealing matches the
minimax-optimal rate for multi-armed bandits with $K$
arms~\citep{auer2002finite}. The EV graph learner is effectively a
bandit over clusters, with the RBM distance providing a geometry on the
arm space: nearby arms have similar payoffs by the EV error bound. This
is analogous to Lipschitz bandits~\citep{kleinberg2005nearly}, where
metric structure enables faster learning. In the EV graph setting, the
metric is not assumed but derived from the game structure.
\end{remark}

\begin{remark}[Tightness]
The exploration term $K V_{\max}$ is tight when each class is first
encountered at a worst-case game with adversarial exploration (zero
payoff). In practice, heuristic strategies achieve reasonable EV, so
the effective exploration cost is much lower. The exploitation term
$(T-K)\varepsilon/2$ is tight at leaves; for trees with branching
factor $m > 1$, Theorem~\ref{thm:ev-error} gives a tighter per-step
bound of $\varepsilon/(2m)$.
\end{remark}

\end{document}
